\section{Popular Attribution Scores as Bilevel Approximations}
\label{sec:attribution-correspondence}

Our first insight is that several apparently different attribution methods can
be placed in one-to-one correspondence with standard ways of approximating the
outer gradient of Eq.~(\ref{eq:bilevel-selection}). The correspondence is made
at the \emph{same} outer state and for the same $f$ and $g$. Canonical data
attribution evaluates the all-present state $\vomega=\bm{1}$; later sections
separately use the budget-feasible state $\vomega^{\mathrm{unif}}$ when the
goal is to optimize a fixed selection budget.

\begin{table}[h]
\centering
\caption{Three attribution families and their corresponding bilevel-gradient
approximations. Each identity is evaluated at a common reference state.}
\label{tab:attribution-correspondence}
\small
\begin{tabularx}{\textwidth}{l l X}
\toprule
Attribution method & Bilevel approximation & Directional identity \\
\midrule
Influence Function & implicit differentiation &
$\bm{s}_{\mathrm{IF}}=-\vg_{\mathrm{IF}}$ \\
TSLOO / SGD-Influence & gradient unrolling &
$\sum_t\bm{s}_{\mathrm{TSLOO}}(t)=-\vg_{\mathrm{GU}}$ \\
Ideal RHO-LOSS & value-function penalty &
$\bm{s}_{\mathrm{RHO}}\propto-\vg_{\mathrm{VF}}$ \\
\bottomrule
\end{tabularx}
\end{table}

\subsection{Implicit differentiation and Influence Functions}

At a lower-level solution $\thetahat(\vomega)$, stationarity gives
$\nabla_{\vtheta}g(\vomega,\thetahat)=0$. Let
\begin{equation}
\rmH(\vomega):=\nabla^2_{\vtheta\vtheta}
g(\vomega,\thetahat(\vomega)),
\qquad
\vq(\vomega):=\nabla_{\vtheta}f(\thetahat(\vomega)).
\end{equation}
When $\rmH$ is invertible, implicit differentiation yields
\begin{equation}
\vg_{\mathrm{IF}}(\vomega)
:=-\nabla^2_{\vomega\vtheta}g(\vomega,\thetahat)
\rmH(\vomega)^{-1}\vq(\vomega).
\label{eq:if-hypergradient}
\end{equation}
For the linear example weights in Eq.~(\ref{eq:fg}), define the validation
Influence Function score \citep{koh2017influence}
\begin{equation}
s_i^{\mathrm{IF}}(\vomega)
:=\frac{1}{n}\nabla_{\vtheta}\ell(\thetahat;z_i)^\intercal
\rmH(\vomega)^{-1}\vq(\vomega).
\label{eq:if-score}
\end{equation}
Concatenating coordinates gives
$\bm{s}_{\mathrm{IF}}(\vomega)=-\vg_{\mathrm{IF}}(\vomega)$. Thus the
classical influence ranking is precisely the descent direction produced by
the implicit-differentiation approximation, evaluated canonically at
$\vomega=\bm{1}$.

\subsection{Gradient unrolling and trajectory-specific influence}

Unrolling replaces the exact lower solution by a finite optimization path,
\begin{equation}
\vtheta_{t+1}(\vomega)=\vtheta_t(\vomega)
-\eta_t\nabla_{\vtheta}g(\vomega,\vtheta_t(\vomega)),
\qquad t=0,\ldots,T-1.
\label{eq:unrolled-training}
\end{equation}
With $\rmH_t=\nabla^2_{\vtheta\vtheta}g(\vomega,\vtheta_t)$,
$\rmG_t=\nabla^2_{\vtheta\vomega}g(\vomega,\vtheta_t)$, and
$\rmJ_t:=\partial\vtheta_t/\partial\vomega^\intercal$, differentiation through
one update gives
\begin{equation}
\rmJ_{t+1}=(\rmI-\eta_t\rmH_t)\rmJ_t-\eta_t\rmG_t,
\qquad
\vg_{\mathrm{GU}}(\vomega)=\rmJ_T^\intercal
\nabla_{\vtheta}f(\vtheta_T).
\label{eq:gu-recurrence}
\end{equation}
For $\rmJ_0=0$, the contribution of example $z_i$ at step $t$ is
\begin{equation}
s_i^{\mathrm{TSLOO}}(t)
:=\frac{\eta_t}{n}\nabla_{\vtheta}\ell(\vtheta_t;z_i)^\intercal
\left[\prod_{k=t+1}^{T-1}(\rmI-\eta_k\rmH_k)\right]
\nabla_{\vtheta}f(\vtheta_T).
\label{eq:tsloo-score}
\end{equation}
This is the first-order trajectory-specific leave-one-out score underlying
SGD-Influence \citep{hara2019sgdinfluence} and subsequent approximate-unrolling
methods \citep{bae2024source,wang2024temporal}. Summing the occurrence-level
scores along the same trajectory gives
\begin{equation}
\sum_{t=0}^{T-1}\bm{s}_{\mathrm{TSLOO}}(t)
=-\vg_{\mathrm{GU}}(\vomega).
\label{eq:tsloo-gu-correspondence}
\end{equation}
Unlike Influence Functions, this direction depends on the realized optimizer,
ordering, and finite training path.

\subsection{Value functions and the branch studied in depth}

The third route replaces the implicit or unrolled hypergradient with a
value-function penalty. Its two-model direction is exactly a per-example loss
difference, which is also the ideal RHO-LOSS score. We derive this third
correspondence in Section~\ref{sec:connection} and then follow it beyond the
canonical one-shot point into multi-round selection.

The unified claim is about \emph{what each score approximates}, not that the
finite estimators are numerically identical. Influence Functions assume a
well-defined local optimum and an inverse Hessian; unrolling targets a realized
finite trajectory; and the value-function route introduces a penalty and two
tracked solutions. All three supply one local outer direction. None, by itself,
solves the constrained outer problem in Eq.~(\ref{eq:bilevel-selection}).
